\chapter{衍生品市场与或有请求权}
\label{chap:Derivatives_Markets_and_Contingent_Claims}

\begin{openingcase}{契约如何在不改变基础业务的情况下改变风险？}

这家制造商计划在三个月内购买1000吨铜。按照当前的价格，这笔采购预计将花费约8000万元人民币。该公司并非试图投机买卖铜价。他们担心的是，如果铜价上涨，生产成本也会相应增加，而在此之前，他们无法及时调整对客户的售价。

该公司可以签订一份远期合约，约定三个月后铜的购买价格。或者，该公司可以选择购买一个看涨期权，该期权可以在铜价大幅上涨时限制成本，同时如果铜价下跌，公司仍能从中受益。无论采取哪种方式，公司对铜的实际需求都不会发生变化。不同的是，公司未来现金流的状态依存模式会有所改变。

如何在不转移基础资产的情况下，让合同重新分配未来不确定情况下的财务后果呢？

\end{openingcase}



\begin{tcolorbox}[
  enhanced,
  colback=FinanceBlue!6,
  colframe=FinanceBlue,
  boxrule=0.6pt,
  sharp corners,
  left=8pt,
  right=8pt,
  top=7pt,
  bottom=7pt,
  before skip=1em,
  after skip=1.6em,
  title=\textbf{核心问题},
  coltitle=white,
  colbacktitle=FinanceBlue,
  fonttitle=\bfseries,
  attach boxed title to top left={xshift=6pt,yshift=-2mm},
  boxed title style={
    boxrule=0pt,
    sharp corners,
    left=5pt,
    right=5pt,
    top=2pt,
    bottom=2pt
  }
]
\begin{enumerate}[
  leftmargin=*,
  label=\textcolor{FinanceBlue}{\textbf{\arabic*.}},
  itemsep=0.55em,
  topsep=0.2em
]


\item 衍生品是如何利用状态依存请求权来拆分和重新分配金融风险敞口的呢？

\item 无套利和复制是如何将衍生品价格与其他交易中的请求权价格联系起来的？

\item 交易所交易、双边合同、抵押品、轧差和中央清算如何影响交易对手风险以及衍生品市场中风险的分配？


\end{enumerate}
\end{tcolorbox}

\section[衍生品存在的原因：或有请求权与风险分配]{衍生品存在的原因：或有请求权与风险分配}

衍生品是支付取决于基础资产、利率、指数、事件或其他指定条件的金融契约。其经济作用不限于投机。衍生品使参与者能够拆分日常经营和金融头寸中捆绑的风险暴露，并将其中选定部分转移给其他市场参与者。

\subsection{状态依存支付}

或有请求权是支付取决于未来状态的金融请求权。设~$T$~时点的可能状态为~$s=1,\ldots,S$，请求权在状态~$s$~下的支付记为~$X_T(s)$。

一个状态代表了与合同支付相关的未来条件。这些条件可能包括基础资产的价格、利率、汇率、商品价格，或者特定信用事件的发生。因此，一项或有请求权描述了合同支付在不同未来状态下的变化方式。

远期契约体现线性的状态依存支付。它约定在未来日期~$T$，按契约订立时确定的交割价格~$K$~买卖基础资产。负有购买义务的一方持有多头，负有交付义务的一方持有空头。

在到期时，长期多头头寸每单位基础资产的支付为

\begin{equation}
X_T^{\mathrm{long}}
=
S_T-K.
\label{eq:ch11-forward-long-payoff}
\end{equation}

短期多头头寸的支付情况则相反：

\begin{equation}
X_T^{\mathrm{short}}
=
K-S_T
=
-X_T^{\mathrm{long}}.
\label{eq:ch11-forward-short-payoff}
\end{equation}

其中，$S_T$~为基础资产到期市场价格。不计交易成本和交易对手损失时，两种头寸的支付等额、方向相反。

期权同样形成状态依存请求权，但持有人拥有行权权利，而非无条件交易义务。其非线性支付结构将在后续章节讨论。

同样的状态依存观点也适用于债务、权益、保险契约和担保。这些请求权在未来的不同状态下分配支付金额，这反映了合同优先权的差异、风险分配的差异，以及支付条件的差异。

\subsection{风险暴露分解与风险拆分}

财务和运营方面的头寸往往同时涉及多种风险来源，这些风险可能会相互作用。例如，一家进口企业的现金流可能受到商品和汇率的影响，而这两个变量之间的共同变动可能会放大或部分抵消该企业本国货币投入成本的变化。因此，一项头寸的风险不仅取决于其各个单独的风险因素，还取决于这些因素之间的相互作用。

风险分解是将整体风险位置分解为经济上不同的来源，并明确每个来源如何影响该位置的现金流或价值。这样就能区分出那些可以保留的风险来源，以及那些可能需要减少或修改的风险来源。

一旦确定了不同的风险敞口，就可以使用衍生品来分别对其进行调整。与特定价格、利率、指数或事件相关的衍生品会创造出一种支付机制，这种支付机制能够响应那些风险来源，而无需改变相关的业务活动或资产状况。通过这种方式，可以将嵌入在更广泛头寸中的风险敞口分离出来，并单独进行调整。

为了说明风险分离，假设在日期$T$时，一个现有头寸的价值为$V_T$，并且面临着多个风险因素。与其中一个因素相关的衍生品支付金额为$X_T$。那么，组合头寸的支付情况就是

\begin{equation}
V_T^{D}=V_T+X_T.
\label{eq:ch11-combined-derivative-position}
\end{equation}

由于~$X_T$~对选定风险因素作出反应，原头寸的这一风险成分可以独立于其他成分调整。如果衍生品支付抵消与该因素相关的部分现金流变化，组合头寸对它的敏感性下降；如果支付强化原有风险暴露，敏感性则可能上升。

套期保值是该机制的一种应用方式。通过采取具有相反支付金额的衍生品头寸，参与者可以在保持原始头寸的其他组成部分基本不变的情况下，减少特定风险敞口。

\subsection{风险分担与市场完备性}

当市场参与者能够交易具有不同状态依存支付结构的请求权时，风险分担变得更加灵活。寻求更稳定现金流的公司可能会将部分风险转移给愿意承担该风险的另一方，以换取补偿，或者因为该风险适合该参与者的其他头寸。衍生品通过创建与特定未来状态相关的合同支付结构，促进了这种风险转移。

风险分担的程度取决于可以从交易中的请求权中构建出的支付范围。考虑一个具有两种未来状态的简单经济体系，分别用$u$和$d$表示。假设有两种交易的证券，它们的支付向量各不相同。

\begin{equation}
\mathbf{x}_1=
\begin{bmatrix}
x_{1u}\\
x_{1d}                                                      
\end{bmatrix},
\qquad
\mathbf{x}_2=
\begin{bmatrix}
x_{2u}\\
x_{2d}
\end{bmatrix}.
\label{eq:ch11-two-security-payoffs}
\end{equation}

持有$a$和$b$两种证券的投资组合，会产生如下状态依存的支付向量：

\begin{equation}
\mathbf{x}_p
=
a\mathbf{x}_1+b\mathbf{x}_2
=
\begin{bmatrix}
ax_{1u}+bx_{2u}\\
ax_{1d}+bx_{2d}
\end{bmatrix},
\label{eq:ch11-spanning-portfolio}
\end{equation}

$a$和$b$的正值表示对相应证券的长期持有，而负值则表示短期抛售。

如果两个支付向量线性独立，它们便张成二维状态支付空间。因此，适当的组合持仓能够产生两种状态下的任意支付向量。所有相关状态依存支付都能由可得交易请求权构造的市场，称为\emph{完备市场}。

如果现有的请求权无法产生所有相关的状态依存支付，那么市场就是不完整的。因此，一些期望的现金流模式就无法精确实现，参与者在交易现有工具后可能仍然承担剩余风险。

在实际市场中，交易成本、交易限制、不可交易的风险、模型的不确定性以及有限的合同类型可能会阻碍某些期望的状态依存支付目标的实现。

衍生品通过引入与现有证券尚未充分代表的风险相联系的请求权，扩展可得的状态依存支付范围，增加参与者之间可转移、可分担的风险种类。衍生品并不创造额外经济资源，而是在不同状态和参与者之间重新分配未来现金流与风险。

\subsection{金融创新与新型请求权}

在不完全市场中，一些状态依存下的期望支付无法从现有的请求权中构建出来。金融创新可以通过创造与特定价格、利率、指数、事件或到期日相关的新衍生品来部分解决这一问题。这些新的请求权可以扩大可转让或对冲的风险范围。

新的请求权可以通过改变参考基准或支付规则来设计。合同可以与新的商品指数、波动性指标或信用事件相关联。其支付结构也可以通过设定上限、下限、障碍特征或路径依赖条件来进行重新设计。通过改变合同的关联因素或支付方式，创新可以改变市场参与者可获得的状态依存现金流。

然而，更紧密的匹配并不能完全消除风险。通过更紧密地跟踪基础资产，定制化的请求权确实可以降低基础风险，但这种做法可能会减少交易的活跃性，或者使得对特定交易对手的依赖程度增加。而标准化的请求权则可以提供更高的流动性，同时与管理中的风险之间的匹配度相对较低。因此，金融创新不仅改变了可以转移的风险类型，也改变了剩余风险的形式。

\section{无套利与复制}

衍生品市场依赖两种相关的机制来保持定价的一致性。无套利机制根据状态依存的支付关系来限制证券的相对价格，而复制机制则利用交易过的证券投资组合来构建等效的支付关系。这两种机制共同为评估衍生品合约提供了基础。

\subsection{套利机会与价格一致性}

一价定律要求，状态依存现金流相同的衍生品头寸具有相同现价。有效比较需要匹配基础参照、名义本金、期限、支付日期、结算币种及其他决定支付的契约条款。这些特征不同可能导致现金流不同，并不必然表明定价不一致。

套利机会可以有两种形式。第一种初始净投资为零，未来所有状态下支付非负，且至少一个状态严格为正。第二种产生正的初始现金流入，却不产生任何未来净支付义务。两种形式均使交易者不承担抵销性未来损失而获得收益。

\paragraph{第一类：从支付优势出发的套利。}

当两种衍生品请求权具有相同的当前价格时，就会出现第一种套利情况。其中一种衍生品在每一个未来状态下提供的价值至少与另一种相同，并且在至少一个状态下其价值严格高于另一种。假设

\begin{equation}
P_A=P_B,
\qquad
X_A(s)\geq X_B(s), \quad \forall s,
\label{eq:ch11-payoff-dominance}
\end{equation}
在至少一种可能的状态下，使用严格的不等式。投资者可以买入$A$，同时卖出$B$。这种策略不需要任何初始净投资：

\begin{equation}
\Pi_0=P_B-P_A=0.
\end{equation}

其未来的净支付满足要求

\begin{equation}
\Pi_T(s)=X_A(s)-X_B(s)\geq 0,
\qquad \forall s,
\end{equation}
在至少一个状态下具有严格正的支付。这种策略能够在不需要初始投资的情况下产生非负的未来支付，从而构成了一种套利机会。

在无套利条件下，如果能够建立相应交易头寸，一项请求权在所有相关状态下支付至少与另一项相等，并在至少一种状态下严格更多，则其现价必须严格更高。

\paragraph{第二种类型：等价支付请求权之间的套利。}

考虑两个交易的衍生品头寸，$A$以及$B$，在每个相关的未来状态下，现金流与当前状态相同。设$P_A$以及那些使用LaTeX命令、公式表达、交叉引用和数字的格式。$P_B$用LaTeX命令、公式、交叉引用和数字来表示它们的当前价格。$X_A(s)$以及那些使用LaTeX命令、公式表达、交叉引用和数字的格式。$X_B(s)$表示它们在状态中的各自支付$s$. 他们的支付等价性
需要

\begin{equation}
X_A(s)=X_B(s),
\qquad \forall s.
\label{eq:ch11-equivalent-derivative-payoffs}
\end{equation}

如果 $P_A < P_B$，那么投资者可以购买资产 $A$，同时做空资产 $B$。初始的净现金流入为……

\begin{equation}
\Pi_0=P_B-P_A>0,
\end{equation}

虽然未来的净支付是

\begin{equation}
\Pi_T(s)=X_A(s)-X_B(s)=0,
\qquad \forall s.
\end{equation}

该策略能够产生正的初始现金流入，同时不会在未来产生净支付义务。如果$P_A>P_B$，那么反向操作这两种头寸就能获得相应的机会。在无套利条件下，相应的衍生品头寸必须具有相同的当前价格。

\begin{example}{支付占优形成的套利}
{ch11-payoff-dominance-arbitrage}

考虑两个具有相同到期日和结算货币的衍生品请求权。它们在两种可能状态下的当前价格和支付情况如下：

\begin{center}
\begin{tabular}{lcc}
\toprule
\textbf{位置}
& \textbf{状态1}
& \textbf{保留LaTeX命令、公式、交叉引用和数字。  
状态2} \\
\midrule
请求权 A（价格 9） & 12 & 20 \\
请求权 B（价格 9） & 10 & 20 \\
\midrule
长文本用LaTeX命令表示，公式用公式编辑器显示，交叉引用则保持不变。数字部分也照原样保留。 & 2 & 0 \\
\bottomrule
\end{tabular}
\end{center}

买入~A、卖空~B~不需要初始净投资，所得头寸在状态~1~产生正支付，在状态~2~支付为零。因此，相同现价不符合无套利条件。在无套利市场中，请求权~A~的现价必须高于请求权~B。

\end{example}

在实际市场中，理论上的套利机会可能无法完全实现。交易成本、买卖价差、卖空限制、融资要求和抵押品需求，以及交易对手风险，都会影响到套利策略的可行性和净收益。 

在没有任何单一证券的支付能够等同于其他支付的情况下，也可以对证券投资组合的支付一致性进行评估。


\subsection{复制与合成头寸}

复制指构造状态依存支付与目标金融请求权相匹配的投资组合。重现另一请求权支付的组合称为\emph{合成头寸}。精确复制时，两种头寸未来现金流完全相同，因此在无套利条件下必须具有相同现值。

考虑一个拥有$J$种交易证券的市场，以及$S$种可能的未来状态。令$X_j(s)$表示在状态$s$下证券$j$的支付额，令$P_j$表示其当前价格。持有每种证券$\theta_j$单位的投资组合会产生相应的支付额。

\begin{equation}
X_p(s)
=
\sum_{j=1}^{J}\theta_j X_j(s),
\qquad s=1,\ldots,S.
\label{eq:ch11-replicating-portfolio-payoff}
\end{equation}

假设一个目标请求权在状态$s$下的支付为$X(s)$。进行复制操作要求投资组合满足

\begin{equation}
\sum_{j=1}^{J}\theta_j X_j(s)
=
X(s),
\qquad s=1,\ldots,S.
\label{eq:ch11-replication-conditions}
\end{equation}

在矩阵表示法中，这个条件可以表示为

\begin{equation}
\mathbf{X}\boldsymbol{\theta}
=
\mathbf{x},
\label{eq:ch11-replication-matrix}
\end{equation}
其中，$\mathbf{X}$ 是交易证券的支付矩阵；$\boldsymbol{\theta}$ 是投资组合持有量的向量；$\mathbf{x}$ 则是目标请求权的支付向量。

当前复制投资组合的成本为

\begin{equation}
P_p
=
\sum_{j=1}^{J}\theta_j P_j.
\label{eq:ch11-replicating-portfolio-cost}
\end{equation}

如果投资组合能够完全复制目标请求权的合同现金流，那么无套利要求该请求权的当前价格$P_X$等于复制成本：

\begin{equation}
P_X=P_p
=
\sum_{j=1}^{J}\theta_j P_j.
\label{eq:ch11-replication-pricing}
\end{equation}

如果目标请求权价格高于复制成本，投资者可以买入复制组合、卖出该请求权；低于复制成本时，则建立相反头寸。两种情况下，未来支付在每一状态相互抵销，投资者获得当前价格差。

\begin{example}{复制或有保障请求权}
{ch11-contingent-protection-replication}

考虑一个拥有三种可能未来状态的经济体：扩张、放缓和严重衰退。市场上交易着三种证券，它们的当前价格分别为人民币约44.20元、人民币约65.00元和人民币约48.00元。这三种证券在三种状态下的每单位支付情况如下所示。

\begin{center}
\begin{tabular}{@{}lccc@{}}
\toprule
\textbf{状态}
& \textbf{证券 1}
& \textbf{证券 2}
& \makecell{\textbf{无风险}\\\textbf{证券}} \\
\midrule
扩展内容       & 100 & 20  & 50 \\
减速现象        & 50  & 50  & 50 \\
严重衰退 & 20  & 100 & 50 \\
\bottomrule
\end{tabular}
\end{center}

假设在扩张期间，商业中断保护请求权的支付为人民币~0；在放缓期间，支付为人民币~50；而在严重衰退期间，支付为人民币~100。其支付向量为

\[
\mathbf{x}
=
\begin{bmatrix}
0\\
50\\
100
\end{bmatrix}.
\]

三种交易证券的投资组合可以再现这种支付情况。让$\theta_1$、$\theta_2$和$\theta_3$分别表示各自的持有量。复制条件为

\begin{align*}
100\theta_1+20\theta_2+50\theta_3&=0,\\
50\theta_1+50\theta_2+50\theta_3&=50,\\
20\theta_1+100\theta_2+50\theta_3&=100.
\end{align*}

解决这些方程可以得到结果。

\[
\theta_1=-0.625,\qquad
\theta_2=0.625,\qquad
\theta_3=1.
\]

因此，该复制投资组合由以下资产组成：0.625单位的证券1的空头头寸，0.625单位的证券2的空头头寸，以及1单位的无风险证券的多头头寸。

其当前成本为

\begin{align*}
P_p
&=-0.625(44.20)+0.625(65)+48\\
&=61.
\end{align*}

该投资组合在所有三种状态下都能实现该保护请求书的支付。在规定的无套利条件下，该请求书的当前价值因此为人民币约61元。

\end{example}

复制不同于更广泛意义上的套期保值。套期保值可能减少现有头寸对特定风险因素的敏感性，却不消除所有状态依存差异；精确复制则要求组合在每一相关状态下都匹配目标请求权的支付。

\section{衍生品市场架构与风险基础设施}

衍生品市场架构决定了交易关系是如何形成的，风险敞口是如何被抵消的，以及交易对手风险是如何被重新分配的。标准化和集中化的交易可以集中流动性并改善价格发现，而抵押品、轧差和清算则限制了双边风险敞口，并重新组织了市场参与者之间损失的分流。

\subsection{交易场所：场内市场与场外市场}

衍生品可以在有组织的交易所进行交易，也可以通过场外交易（OTC）的方式交易。在交易所交易的衍生品通常采用标准化的合约条款，包括合约规模、到期日、结算条款以及其他规格。标准化使得许多参与者能够使用具有共同条款的合约进行交易，从而使各种头寸能够更方便地互换。

场外衍生品交易是由交易对手之间协商达成的，通常遵循标准化的法律文件，同时允许根据经济条件进行定制。因此，名义金额、到期日、参考变量和结算条款都可以进行调整，以更紧密地匹配特定的风险敞口。然而，更大的定制化可能会使得合约的互换性降低，从而增加了头寸转移的难度。

这一区别并非绝对。有些场外衍生品条款高度标准化，并可能进行中央清算；电子交易平台也缩小了传统上场内执行与交易商执行的差异。因此，交易场所、契约设计与清算安排是相关但不同的市场组织维度。

主要的衍生品交易所提供了标准化合约如何在正式交易场所中组织的具体例子。表~\ref{tab:ch11-major-derivative-exchanges}展示了来自中国和主要国际市场的精选案例。

\begin{table}[H]
\centering
\caption{主要场内衍生品市场示例}
\label{tab:ch11-major-derivative-exchanges}

\begin{tabularx}{\textwidth}{p{3.8cm} p{3.2cm} X}
\toprule
\textbf{交换} &
\textbf{市场焦点} &
\textbf{代表性衍生品} \\
\midrule

中国金融期货交易所 (CFFEX) & 衍生品 & 权益指数期货与期权；政府债券期货 \\

上海期货交易所 (SHFE) &  
商品衍生品 &  
金属、能源及其他商品期货和期权 \\

大连商品交易所 (DCE) & 
衍生品 & 
农业、工业及化工期货与期权 \\

郑州商品交易所（ZCE） &
衍生品市场 &
农业、能源及工业期货与期权 \\

广州期货交易所 (GFEX) &  
商品及新兴衍生品 &  
与特定行业和新兴市场相关的合约 \\

上海和深圳证券交易所 & 
证券衍生品 & 
交易所交易ETF及其他批准的期权 \\

\addlinespace
CME小组 &
金融和商品衍生品 &
利率衍生品、权益指数衍生品、外汇衍生品和商品衍生品 \\

Eurex &
金融衍生品 &
权益指数、利率及其他衍生品 \\

ICE &
金融和商品衍生品 &
能源、利率、权益及其他衍生品 \\

\bottomrule
\end{tabularx}
\end{table}

\subsection{流动性与价格发现}

衍生品市场的设计会影响交易活动和信息的集中程度。标准化合约可以将具有相似风险敞口的参与者聚集在同一工具上，这可能会提高交易的流动性，并使市场价格更加容易观察。而定制化的合约则能够更紧密地匹配各个个体的风险敞口，但交易往往分散在更多不同的头寸上。

液体市场也有助于价格的确定。当参与程度较高且交易频繁发生时，报价和交易价格就能吸收新的信息，为相关工具和基础市场提供参考价值。这种信息功能取决于市场的深度、参与度和交易条件；仅有标准化是不够的，它无法确保市场具有活跃性和信息性。


\subsection{交易对手风险暴露、抵押品与轧差}

衍生品的市场价值会随着市场状况的变化而发生变化。当双边合约对某一方具有正面价值时，该方就对交易对手方拥有请求权，如果交易对手方未能履行义务，该方可能会遭受损失。因此，交易对手方的风险敞口会随着合约的市场价值而变化，并且会随时间推移而发展。

交易对手信用风险不同于衍生品市场风险。市场风险来自基础价格、利率或其他风险因素变化，交易对手信用风险则来自契约义务可能无法履行。当市场变化增加风险暴露、同时交易对手信用质量恶化时，两种风险可能相互作用。

保证金要求提供了随衍生品价值变化控制交易对手风险暴露的机制。两类保证金作用不同：初始保证金是覆盖违约头寸平仓或替代所需期间潜在未来风险暴露的抵押品；变动保证金反映契约当前市场价值变化，并随价值变化在交易对手之间转移。

在交易型期货市场中，合约在每个交易日结束时根据市场情况进行结算。每日的盈亏是根据结算价格的变动来计算的，并计入或从参与者的保证金账户中扣除。对于涉及$n$份合约且合约乘数為$m$的长期期货头寸，每日的保证金变化为

\begin{equation}
VM_t=nm(F_t-F_{t-1}),
\label{eq:ch11-daily-variation-margin}
\end{equation}

其中，$F_t$和$F_{t-1}$分别表示当前和之前的结算价格。持有短期头寸的人需要相反的变动保证金。

每日结算在损益发生时转移损益，而非让市场价值变化积累为无担保契约风险暴露。初始保证金则为参与者违约后平仓期间的潜在损失提供额外保护。

抵押品与轧差是降低双边无担保风险暴露的两项重要机制。抵押品随契约价值变化，为义务提供资产支持。具有法律可执行性的轧差，允许合并考虑同一交易对手之间的正负头寸，而非逐笔分别处理。轧差减少可能需要抵押品支持的风险暴露，抵押品则减少剩余义务中的无担保部分。

\subsection{中央清算与交易对手风险再分配}

抵押品和轧差降低双边关系中的交易对手风险暴露，中央清算则改变这些关系的结构。合格衍生品被接受清算后，中央对手方（CCP）介入原交易双方之间，各清算成员就已清算义务面对~CCP。原本双边存在的风险暴露因此在清算体系中重新组织。

\begin{figure}[htbp]

\centering

\begin{tikzpicture}[
    >=Latex,
    dealer/.style={
        draw,
        rounded corners,
        minimum width=20mm,
        minimum height=8mm,
        align=center,
        font=\small
    },
    ccp/.style={
        draw=FinanceBlue,
        fill=FinanceBlue!8,
        rounded corners,
        minimum width=24mm,
        minimum height=12mm,
        align=center,
        font=\small
    },
    bilateral/.style={
        Latex-Latex,
        line width=0.7pt,
        draw=FinanceBlue
    },
    title/.style={
        font=\small\itshape
    }
]

% Left panel: bilateral structure
\node[title] at (2.5,3.5) {双边市场结构};

\node[dealer] (A1) at (0.8,2.4) {交易商A};
\node[dealer] (B1) at (4.2,2.4) {交易商B};
\node[dealer] (C1) at (0.8,0.4) {交易商 C};
\node[dealer] (D1) at (4.2,0.4) {交易商 D};

\draw[bilateral] (A1) -- (B1);
\draw[bilateral] (A1) -- (C1);
\draw[bilateral] (A1) -- (D1);
\draw[bilateral] (B1) -- (C1);
\draw[bilateral] (B1) -- (D1);
\draw[bilateral] (C1) -- (D1);

\node[align=center, font=\scriptsize]
at (2.5,-0.55)
{多次双边曝光};

% Arrow between structures
\draw[->, line width=0.9pt, gray!70]
(5.25,1.4) -- (6.25,1.4);

% Right panel: centrally cleared structure
\node[title] at (9.0,3.5) {中央统一整理的结构};

\node[dealer] (A2) at (7.3,2.5) {交易商A};
\node[dealer] (B2) at (10.7,2.5) {交易商B};
\node[dealer] (C2) at (7.3,0.3) {交易商 C};
\node[dealer] (D2) at (10.7,0.3) {交易商 D};

\node[ccp] (CCP) at (9.0,1.4) {中心\\交易对手方};

\draw[line width=0.8pt, FinanceBlue] (A2) -- (CCP);
\draw[line width=0.8pt, FinanceBlue] (B2) -- (CCP);
\draw[line width=0.8pt, FinanceBlue] (C2) -- (CCP);
\draw[line width=0.8pt, FinanceBlue] (D2) -- (CCP);

\node[align=center, font=\scriptsize]
at (9.0,-0.55)
{通过CCP重新组织了各部分的布局。};

\end{tikzpicture}

\caption{双边与中央清算下的交易对手关系}

\label{fig:ch11-bilateral-central-clearing}

\end{figure}

中央清算还可以支持跨清算头寸的多边轧差。那些原本需要分别处理的两岸关系中的义务，可以在清算系统中进行抵消，从而减少某些总风险敞口。

因此，中央清算重新分配交易对手风险，而非消除风险。交易对手关系围绕~CCP~组织，而不是分散在许多双边关系中，使~CCP~成为已清算市场的核心节点。

%========================================================
% Summary
%========================================================

\TailModule{summary}{本章小结}

衍生品市场为交易合同提供了制度框架，这些合同的现金流取决于基础资产、利率、指数或特定事件。这些合同产生了或有请求权，使得支付可以在未来的不同状态下进行分配。通过将特定的风险暴露与更广泛的商业和财务头寸分离开来，衍生品有助于市场参与者之间的风险转移和分担。可获得的或有请求权的范围取决于可交易的金融工具以及市场的完整性。

或有请求权的价格通过交易头寸之间的无套利关系相互关联。支付的主导性对相对价格施加了限制，而具有相同状态依存支付的请求权必须具有相同的当前价格。当一项请求权可以通过一组交易证券的投资组合来复制时，其价格必须等于该投资组合的成本。因此，交易工具的可用性决定了哪些请求权可以精确复制，而对于那些无法复制的请求权，则存在一系列符合无套利原则的支付价格范围。

衍生品市场架构决定了交易、结算以及交易对手风险管理的模式。交易所交易市场和场外交易市场在标准化程度及合同灵活性上存在差异。抵押品、交易对手风险管理等因素用于管理交易对手风险，而清算机构则通过中央对手方来重新安排双边交易的风险。这些安排虽然能够重新分配风险，但并不能完全消除风险。

%========================================================
% Key Terms
%========================================================

\TailModule{keyterms}{关键术语}
或有请求权（Contingent claims）；状态依存支付（state-contingent payoffs）；风险暴露分解（exposure decomposition）；风险拆分（risk unbundling）；风险分担（risk sharing）；市场完备性（market completeness）；金融创新（financial innovation）；一价定律（law of one price）；套利（arbitrage）；支付占优（payoff dominance）；复制（replication）；合成头寸（synthetic positions）；场内衍生品（exchange-traded derivatives）；场外衍生品（over-the-counter (OTC) derivatives）；流动性（liquidity）；价格发现（price discovery）；交易对手风险暴露（counterparty exposure）；保证金要求（margin requirements）；抵押品与轧差（collateral and netting）；中央清算与中央对手方（central clearing and central counterparty (CCP)）。

%========================================================
% Financial AI Lab
%========================================================

\TailModule{aiprompt}{金融~AI~实验室}

%========================================================
% Financial AI Lab
%========================================================

\TailModule{aiprompt}{金融~AI~实验室}

%========================================================
% Financial AI Lab
%========================================================

\TailModule{aiprompt}{金融~AI~实验室}

\textbf{实验：在衍生品市场中审计交易对手风险与风险管理}

请一个AI系统来评估一家公司OTC衍生品投资组合中的交易对手风险敞口。假设有一家公司在同一日期拥有以下头寸，所有金额均以人民币百万为单位计算。

\begin{center}
\begin{tabular}{lcc}
\toprule
\textbf{对方银行}
& \textbf{合同/协议}
& \textbf{市场价值} \\
\midrule
交易商A & 合同1 & $+6.0$ \\
交易商A & 合同2 & $-4.0$ \\
交易商B & 合同3 & $+3.0$ \\
\bottomrule
\end{tabular}
\end{center}

与交易商A签订的 두 계약은 동일한 법적으로 구속력 있는 종료轧差 협정에 따릅니다. 거래商A는 적용 가능한抵押品折扣率에 따라 회사에 약 120만 RMB의 변동 마진을 지불했습니다. 반면, 거래商B와의 계약은 별도의轧差 협정에 따르며, 아직 어떠한抵押品도 받지 않았습니다.

要求人工智能系统计算该公司当前未受担保的对手方风险敞口，并评估不同的风险管理安排对该风险敞口的影响。

将答案分为四个阶段进行审核。

\begin{enumerate}

\item \textbf{计算交易对手的净风险敞口。}

确定在应用具有法律约束力的轧差后，对每个交易商的净当前敞口。考虑所收到的抵押品，并解释为什么与不同交易商持有的头寸不能自动相互抵消。

\item \textbf{区分市场风险与交易对手信用风险。}

解释基础价格的变化如何影响合约的市场价值，以及公司对交易对手的潜在风险。阐明为什么正的置换价值会使得公司面临交易对手无法履行义务的风险。

\item \textbf{评估保证金和抵押品安排。}

请保留LaTeX命令、公式、交叉引用和数字的格式不变。

区分初始保证金与变动保证金的区别。
解释每日市价跟踪机制和变动保证金机制如何影响当前的持仓状况，同时说明初始保证金如何能够在清算期间保护投资者免受未来潜在风险的影响。指出即使进行了抵押担保，仍可能存在的风险。

\item \textbf{评估中央清算的作用。}

解释引入中央交易对手如何改变原有的双边关系。
讨论多边轧差、保证金要求、
违约风险以及交易对手风险在清算系统中的集中情况。
区分交易对手风险的重新分配与市场风险消除的区别。

\end{enumerate}

最终报告应验证抵押品的计算结果，明确与轧差和抵押品相关的法律和运营假设，并区分风险降低与风险转移以及风险集中。

%========================================================
% Discussion Questions
%========================================================

\TailModule{questions}{讨论题}

\begin{enumerate}

\item 衍生品如何在不改变公司基础业务活动的情况下，对选定的金融风险进行调节？

\item 市场完整性如何影响复制或有请求权的能力？在无法完全复制的情况下，可能会存在哪些定价不确定性？

\item 为什么定制化的OTC合约能够在增加抵押品、保证金要求以及双边轧差的重要性同时，使得对基础资产的风险敞口更加匹配呢？

\item 中央清算是如何重新安排交易对手关系的？为什么这会使中央交易对手的风险管理和韧性变得重要？

\end{enumerate}

%========================================================
% Applied Problems
%========================================================

\TailModule{appprobs}{应用题}

\begin{enumerate}

\item \textbf{状态依存的支付和远期套期保值。}

一位商品生产商预计在日期$T$出售一单位该商品。在该日期的现货价格可能为70、100或140。该生产商签订了一份现金结算的短期合约，合约的交付价格为100，到期日亦然。

\begin{enumerate}

\item 计算每个状态下短期远期头寸的支付情况。

\item 计算生产者在每个状态下的物理销售和未来头寸所带来的现金流总和。

\item 解释远期合约如何改变生产者的现金流在不同州的分配情况。

\end{enumerate}


\item \textbf{暴露分解与风险分离。}

该公司在日期$T$时将购买500个单位的进口商品。以外币计价的商品价格表示为$P_T$，而单位外币在本国货币中的价格则表示为$S_T$。该公司使用本国货币进行购买时的成本为……

\[
C_T=500P_TS_T.
\]

\begin{enumerate}

\item 确定影响企业购买成本的两大市场风险因素。

\item 解释每个风险因素的变化如何影响公司的本国货币现金流出。

\item 假设该公司使用的衍生品仅与商品价格$P_T$挂钩。哪些风险敞口会发生变化，哪些则保持不变？

\end{enumerate}


\item \textbf{或有请求权的复制。}

在双状态市场中交易了两份证券。它们当前的价格和支付情况如下：

\begin{center}
\begin{tabular}{lccc}
\toprule
\textbf{证券}
& \textbf{当前价格}
& \textbf{上状态}
& \textbf{关闭状态} \\
\midrule
证券 A & 8 & 12 & 6 \\
证券 B & 9 & 4 & 10 \\
\bottomrule
\end{tabular}
\end{center}

在“上状态”下，或有请求权可获得20的赔偿；在“下状态”下，可获得2的赔偿。

\begin{enumerate}

\item 确定需要持有的证券A和B的数量，以履行该或有请求权。

\item 请标明哪些部分是长句，哪些部分是短句。同时保留所有LaTeX命令、公式表达、交叉引用以及数字的内容。

\item 计算当前投资组合的成本。

\item 解释为什么在无套利条件下，请求权的当前价值必须等于复制成本。

\end{enumerate}


\item \textbf{套利机会与价格一致性。}

考虑两对衍生品请求权。

对于第一对请求权，请求权A和B在状态1时的支付均为15，在状态2时的支付均为8。它们当前的价值分别为13和10。

对于第二对请求权，C和D的当前价格都是11。它们的支付情况如下：

\begin{center}
\begin{tabular}{lcc}
\toprule
\textbf{请求权}
& \textbf{状态1}
& \textbf{保留LaTeX命令、公式、交叉引用和数字。  
状态2} \\
\midrule
请求权 C & 14 & 20 \\
请求权 D & 12 & 20 \\
\bottomrule
\end{tabular}
\end{center}

\begin{enumerate}

\item 对于第一对情况，确定套利策略，并计算其初始现金流以及未来的支付金额。

\item 对于第二对情况，确定套利策略，并计算其初始现金流以及未来的支付金额。

\item 为什么预期回报较高的投资不一定构成套利机会呢？

\end{enumerate}


\end{enumerate}
%========================================================
% Further Reading
%========================================================

\TailModule{furtherread}{延伸阅读}

\begin{itemize}

\item Duffie, Darrell, and Chi-Fu Huang (1985),
``Implementing Arrow-Debreu Equilibria by Continuous Trading of Few Long-Lived Securities（通过少数长期证券的连续交易实现阿罗—德布鲁均衡）,''
\emph{Econometrica}, 53(6), 1337--1356.

\end{itemize}

\ChapterTail{summary,keyterms,aiprompt,questions,appprobs,furtherread}
